\chapter{Application Streamlit et déploiement}

% Partie consacrée à l'application utilisateur et au bonus de déploiement.

\section{Interface utilisateur}
% - Présentation de \texttt{streamlit\_app.py}.
% - Fonctionnalités :
%   \begin{itemize}
%     \item zone de question,
%     \item affichage de la réponse générée,
%     \item affichage des documents sources avec scores,
%     \item visualisation du reranking (score de similarité + overlap),
%     \item gestion du cas « pas de réponse fiable ».
%   \end{itemize}
% - Captures d'écran à insérer.

%\begin{figure}[ht]
%    \centering
%    \includegraphics[width=0.85\textwidth]{assets/screenshot_streamlit.png}
%    \caption{Capture d'écran de l'application Streamlit.}
%    \label{fig:streamlit}
%\end{figure}

\section{Configuration et reproductibilité}
% - Variables d'environnement (\texttt{.env.example}).
% - Installation, ingestion, lancement.
% - Choix de l'API UTC pour le LLM et les embeddings.

\section{Déploiement (bonus)}
% - Mode de déploiement choisi (Streamlit Community Cloud, Docker, serveur
%   UTC\dots).
% - Difficultés rencontrées (gestion des clés, persistance Chroma, taille
%   du corpus, etc.).
% - Lien vers l'application déployée si applicable.
